\chapter{五类特征转换方法与关键算法}
\label{chap:4}
\chapterguide{以“问题、输入、步骤、输出、适用条件”组织方法，减少接口名称。关键公式和案例应与真实代码路径及可复现实验对应。}

\section{建模引用的恢复方法}
\label{sec:4-1}

\subsection{源目标对象对应与引用解析}
\label{subsec:4-1-1}
\outlinehint{说明源标识到统一对象、统一对象到目标原生对象的对应，依据历史顺序使用当前已创建的引用。}

\subsection{基于几何信息的目标对象匹配}
\label{subsec:4-1-2}
\outlinehint{利用类型、位置、方向、采样点和必要上下文定位目标边面，处理分段和几何归属变化。}

\subsection{多候选、失效引用与失败处理}
\label{subsec:4-1-3}
\outlinehint{讨论重复或对称几何产生的歧义，明确可用判定、拒绝条件与日志，不把全部候选歧义写成已解决。}

\section{基准面的关联重建方法}
\label{sec:4-2}

\subsection{构造约束与参考组合映射}
\label{subsec:4-2-1}
\outlinehint{按建面方法和参考角色组织目标输入，说明过点、过线、偏置和角度条件的组合与必要性。}

\subsection{偏置方向、角度分支与建后核验}
\label{subsec:4-2-2}
\outlinehint{依据源平面位置和方向选择目标分支，创建后核对实际平面；使用固定辅助点时单独记录关联变化。}

\section{草图几何与约束的恢复方法}
\label{sec:4-3}

\subsection{源目标草图坐标转换}
\label{subsec:4-3-1}
\outlinehint{用源局部坐标到世界坐标、再到目标实际局部坐标的过程推导变换，说明坐标系正交性和单位前提。}

\subsection{圆弧区间与构造几何的恢复}
\label{subsec:4-3-2}
\outlinehint{结合圆心、端点、方向和必要中间点确定同一物理弧段，保留中心线和辅助图元的作用。}

\subsection{约束作用位置与关联恢复}
\label{subsec:4-3-3}
\outlinehint{映射约束参与对象和端点等位置，校正水平、竖直等方向关系；区分驱动尺寸恢复与仅几何复制。}

\section{拉伸方向与终止范围的转换方法}
\label{sec:4-4}

\subsection{实际方向与材料操作的统一}
\label{subsec:4-4-1}
\outlinehint{将源端方向标记转换为明确的有效方向，再映射目标加料和切除路径，避免重复翻转。}

\subsection{单向、双向与对称深度转换}
\label{subsec:4-4-2}
\outlinehint{用两侧材料范围解释深度，对对称总量与单侧量给出换算，并保留双端各自的终止规则。}

\subsection{引用终止与终止偏移处理}
\label{subsec:4-4-3}
\outlinehint{恢复终止对象与偏移，区分有限面和延伸支承面；明确辅助构造、回退定深及拒绝处理对关联的影响。}

\section{旋转方向与角度范围的转换方法}
\label{sec:4-5}

\subsection{旋转轴定位与正转方向判定}
\label{subsec:4-5-1}
\outlinehint{分别恢复轴线位置和公共正向；说明Creo中结合有效轮廓、轴线与草图法向进行判定的依据和退化边界。}

\subsection{双向、对称与正负角度转换}
\label{subsec:4-5-2}
\outlinehint{以实际扫过角域解释两侧参数，使用正90度、反30度对应负30度到正90度的范围作为示例。}

\subsection{参考终止与互补角问题处理}
\label{subsec:4-5-3}
\outlinehint{说明到面、到点的辅助对象与取解要求，明确互补角的未完成范围以及共享偏移量纲的专项检查。}

\section{薄壁材料位置的统一映射方法}
\label{sec:4-6}

\subsection{两侧偏置与材料区间的对应}
\label{subsec:4-6-1}
\outlinehint{用轮廓两侧边界位置表达薄壁，解释单侧、对称及不等量双侧的总厚度和位置差异。}

\subsection{三系统厚度选项映射与限制}
\label{subsec:4-6-2}
\outlinehint{归纳厚度加方向、材料侧与有符号偏置间的映射；明确不等量双侧信息损失等现有边界。}

\section{扫掠截面与路径的联合重建方法}
\label{sec:4-7}

\subsection{路径起点、曲线顺序与方向恢复}
\label{subsec:4-7-1}
\outlinehint{按端点连续性和源选取信息组织路径，说明路径反向、闭合路径与跨对象曲线链的不同处理范围。}

\subsection{截面放置与目标坐标适配}
\label{subsec:4-7-2}
\outlinehint{联合路径起点、切向和截面实际坐标恢复放置，说明截面几何不变仍可能因方向不同生成不同结果。}

\subsection{独立、内嵌与参数圆截面的重建}
\label{subsec:4-7-3}
\outlinehint{分别说明直接引用、补建草图、截面回放以及圆管或双圆截面构造，并检查内半径是否参与生成。}

\subsection{特殊截面补偿与重建策略边界}
\label{subsec:4-7-4}
\outlinehint{解释斜交截面的受限投影补偿及适用条件，区分坐标重表达与真正几何投影；记录重试时合并或方向选项的变化。}

\section{本章小结}
\label{sec:4-8}

\outlinehint{归纳本章能够支持的认识，并说明与下一章的衔接；不重复列举所有接口或扩展未证明结论。}
